% !TEX program = xelatex
% Compile with: xelatex SaundersAdam_CV.tex && biber SaundersAdam_CV && xelatex SaundersAdam_CV.tex && xelatex SaundersAdam_CV.tex
\documentclass[11pt]{article}

\usepackage[letterpaper, margin=1in]{geometry}
\usepackage{fontspec}
\setmainfont{Times New Roman}
\usepackage{enumitem}
\usepackage{xcolor}
\usepackage{etoolbox}

\setlength{\parindent}{0pt}
\setlength{\parskip}{0pt}
\emergencystretch=1.5em
\pagestyle{empty}

% ---------------------------------------------------------------------------
% Layout machinery: a fixed-width "label" column on the left and a "content"
% column on the right. The label is printed only on the first row of a
% section; every following row/paragraph passes an empty label so the text
% keeps flowing at the content-column indent -- including across page breaks,
% since this is just ordinary paragraphs, not a table. Labels are kept to a
% single line (the column is sized to fit the widest one, "Abstracts/Short
% Papers") so a label never spills onto a second line and leaves a blank gap
% next to it.
% ---------------------------------------------------------------------------
\newlength{\cvlabelwidth}
\setlength{\cvlabelwidth}{1.2in}
\setlist[itemize]{leftmargin=\dimexpr\cvlabelwidth+0.25in\relax,
                  label=\textbullet, itemsep=1pt, topsep=3pt, parsep=0pt,
                  partopsep=0pt}

% \cvfitbox{width}{content} -- a fixed-width, top-aligned box that lets
% its content wrap onto further lines inside it instead of overflowing --
% but a top-aligned multi-line \parbox reports its *entire* extra height
% as depth (everything below its first baseline), so without correction
% the very next line of running text after it gets pushed down by all of
% that accumulated depth in one lump, i.e. one visible gap exactly where
% a wrapped box meets whatever comes after it. Forcing the depth back
% down to a normal single line's depth (via a reference \strut) doesn't
% touch what's drawn -- the box's later lines still print in full -- it
% only corrects what TeX *thinks* the box's footprint is, so the next
% line lands at the usual single-line spacing instead. \strut inside the
% box also guards against a real oddity: a *completely* empty \parbox
% (an empty label or date) silently reverts to zero width, which would
% throw off that row's own indent/alignment.
\newbox\cvfitboxreg
\newbox\cvrefboxreg
\newcommand{\cvfitbox}[2]{%
  \setbox\cvfitboxreg=\hbox{\parbox[t]{#1}{\raggedright#2\strut}}%
  \setbox\cvrefboxreg=\hbox{\strut}%
  \dp\cvfitboxreg=\dp\cvrefboxreg\relax
  \box\cvfitboxreg}

% \cvlabelbox{content} -- \cvfitbox at the fixed label-column width, for
% use inline at the start of a \hangindent paragraph. A label too wide
% for \cvlabelwidth (e.g. "Scientific Experience") wraps onto a second
% line inside it, flush with the content that follows thanks to
% \cvfitbox's depth correction above.
\newcommand{\cvlabelbox}[1]{\cvfitbox{\cvlabelwidth}{#1}}

% \cvplainbox{width}{content} -- like \cvfitbox, but *without* the depth
% correction: for \cvrow's content box below. A wrapped label is a
% section heading whose leftover line(s) are meant to be shared with
% whatever short entry comes right after (that's what \cvfitbox's
% correction is for), but a wrapped *entry* -- an award name, say -- is
% that entry's own real text, and the row after it must not start until
% that text is actually done printing. So the content box keeps its true
% (uncorrected) depth, which is exactly the extra space needed to clear
% it, while the label box beside it still uses the corrected \cvlabelbox.
% \hangindent inside the box indents every line after the first, so a
% wrapped second (or third) line reads as a clear continuation of the
% first rather than lining back up with it.
\newcommand{\cvplainbox}[2]{\hbox{\parbox[t]{#1}{\hangindent=1em\hangafter=1\raggedright#2\strut}}}

% \cvrow{label}{text}{date} -- entry row with a right-flushed date. The
% date column is a fixed width reserved out of the line, so long entry
% text wraps onto a further line -- inside its own \cvplainbox, narrower
% than the full remaining width -- rather than running up against (or
% past) the date. Whichever of label, content is taller wins: a wrapped
% label's extra lines are free to be shared with a short next entry (see
% \cvplainbox above), but a wrapped entry's own extra lines are real,
% reserved space. \cvdatewidth is sized generously for the longest date
% range used in this CV (e.g. "August 2019 -- May 2023"); \makebox keeps
% the date text together as a unit so it can't split mid-phrase, e.g.
% "May" / "2023".
\newlength{\cvdatewidth}
\setlength{\cvdatewidth}{1.5in}
\newlength{\cvcontentwidth}
\setlength{\cvcontentwidth}{\dimexpr\textwidth-\cvlabelwidth-\cvdatewidth-0.5em\relax}
\newcommand{\cvrow}[3]{%
  \par\noindent
  \cvlabelbox{\textbf{#1}}\cvplainbox{\cvcontentwidth}{#2}\hspace{0.5em}\makebox[\cvdatewidth][r]{#3}\par}

% \cvpar{label}{text} -- plain running paragraph, no date column
\newcommand{\cvpar}[2]{%
  \par\noindent\hangindent=\cvlabelwidth\hangafter=1
  \cvlabelbox{\textbf{#1}}#2\par}

% Inline red highlight (matches the award callout in the source CV)
\newcommand{\highlight}[1]{\textcolor{red}{#1}}

% ---------------------------------------------------------------------------
% Publications, pulled from separate BibTeX files by status. See cvbib.bbx
% for the custom bibliography driver (author list with the CV owner's name
% auto-bolded, plus the "cvmain"/"cvtag" layout environments used below).
% ---------------------------------------------------------------------------
% sorting=ydnt sorts each \printbibliography call (already split into
% accepted/published/submitted groups by its own keyword filter) by year,
% descending -- newest first, matching the hand-ordered \cvrow sections.
\usepackage[backend=biber, sorting=ydnt, bibstyle=cvbib, citestyle=numeric, maxnames=99, minnames=99]{biblatex}
% citestyle=numeric's own .cbx loads after cvbib.bbx and re-declares the
% title format with quotation marks, so override it here instead, after
% both have loaded -- cvbib.bbx's driver adds its own quotation marks
% around the title explicitly, so the field format itself must stay
% plain. An entrytype-specific format takes priority over the generic
% one, so override article and inproceedings (conference papers,
% abstracts, talks) too.
\DeclareFieldFormat{title}{#1}
\DeclareFieldFormat[article]{title}{#1}
\DeclareFieldFormat[inproceedings]{title}{#1}
% biblatex's standard style italicizes both journaltitle and booktitle by
% default (each wrapped in \mkbibemph, i.e. \emph). journaltitle's default
% italic is exactly what's wanted for a journal name, so it's left alone;
% booktitle's is reset to plain here so a conference/workshop name (used
% for conference papers, abstracts, and talks) prints upright instead.
\DeclareFieldFormat[inproceedings]{booktitle}{#1}
% Likewise, citestyle=numeric's default doi format already prepends its
% own "DOI:" label (and wraps it as a link), so it must be reset to a bare
% value too -- cvbib.bbx's driver supplies the "doi:" label itself.
\DeclareFieldFormat{doi}{#1}
\addbibresource[label=journal-published]{bib/journal-published.bib}
\addbibresource[label=journal-accepted]{bib/journal-accepted.bib}
\addbibresource[label=journal-submitted]{bib/journal-submitted.bib}
\addbibresource[label=conference-published]{bib/conference-published.bib}
\addbibresource[label=conference-accepted]{bib/conference-accepted.bib}
\addbibresource[label=conference-submitted]{bib/conference-submitted.bib}
\addbibresource[label=talks]{bib/talks.bib}
\addbibresource[label=abstracts]{bib/abstracts.bib}
\nocite{*}

\begin{document}

% ------------------------------- Header -------------------------------
\begin{center}
  {\LARGE\textbf{Adam Saunders}}
\end{center}
\noindent\rule{\linewidth}{0.4pt}

\noindent\textbf{Email:} adam.m.saunders@vanderbilt.edu \hfill \textbf{Phone:} (937) 620-5338

\vspace{10pt}

% ------------------------------- Summary -------------------------------
\cvpar{Summary}{An electrical engineering PhD student studying medical imaging. Research focused on applying artificial intelligence tools to impute missing information from imperfect data. Strengths in analytical skills and scientific communication.}

\vspace{10pt}

% ------------------------------- Education -------------------------------
\cvrow{Education}{PhD, Electrical and Computer Engineering}{July 2023 -- Present}
\cvrow{}{\textit{Vanderbilt University}}{}
\cvrow{}{Bachelor of Electrical Engineering,\hspace{0.3em}\textit{summa cum laude}}{August 2019 -- May 2023}
\cvrow{}{\textit{University of Dayton}}{}
\begin{itemize}
  \item Minor in mathematics
\end{itemize}

\vspace{10pt}

% ------------------------------- Scientific Experience -------------------------------
\cvrow{Scientific Experience}{Research Assistant}{July 2023 -- Present}
\cvrow{}{\textit{Vanderbilt University}}{}
\begin{itemize}
  \item Researcher in the Medical-image Analysis and Statistical Interpretation (MASI) Lab and Vanderbilt Lab for Immersive AI Translation (VALIANT)
  \item Collaborated with scientists to communicate published findings
\end{itemize}

\cvrow{}{SULI Intern}{June 2022 -- August 2022}
\cvrow{}{\textit{Oak Ridge National Laboratory}}{}
\begin{itemize}
  \item Explored deep learning algorithms on a world-leading supercomputer
  \item Performed and presented medical image processing research
\end{itemize}

\cvrow{}{Physics Teaching Assistant}{January 2022 -- May 2022}
\cvrow{}{\textit{University of Dayton}}{January 2021 -- May 2021}
\begin{itemize}
  \item Empowered over 50 students to learn calculus-based mechanics
  \item Assisted in the development of testing materials and led office hours
\end{itemize}

\cvrow{}{Computer Engineering Intern}{May 2021 -- August 2021}
\cvrow{}{\textit{Defense Research Associates, Inc.}}{}
\begin{itemize}
  \item Designed and implemented digital hardware interfaces for sensors
  \item Presented technical results and documentation for future deployment
\end{itemize}

\cvrow{}{Physics Grader}{January 2020 -- May 2020}
\cvrow{}{\textit{University of Dayton}}{}
\begin{itemize}
  \item Evaluated exams and group exercises for over 50 students
  \item Led weekly homework help sessions in a challenging, fast-paced class
\end{itemize}

\vspace{10pt}

% ------------------------------- Journal Articles -------------------------------
\global\cvbibfirsttrue
\renewcommand{\cvbiblabel}{Journal Articles}
\renewcommand{\cvbibtag}{accepted}
\printbibliography[keyword=journal-accepted, env=cvpub, heading=none]

\vspace{10pt}
\renewcommand{\cvbibtag}{}
\printbibliography[keyword=journal-published, env=cvpub, heading=none]

\vspace{10pt}
\renewcommand{\cvbibtag}{submitted}
\printbibliography[keyword=journal-submitted, env=cvpub, heading=none]

\vspace{10pt}

% ------------------------------- Conference Papers -------------------------------
\global\cvbibfirsttrue
\renewcommand{\cvbiblabel}{Conference Papers}
\renewcommand{\cvbibtag}{accepted}
\printbibliography[keyword=conference-accepted, env=cvpub, heading=none]

\vspace{10pt}
\renewcommand{\cvbibtag}{}
\printbibliography[keyword=conference-published, env=cvpub, heading=none]

\vspace{10pt}
\renewcommand{\cvbibtag}{submitted}
\printbibliography[keyword=conference-submitted, env=cvpub, heading=none]

\vspace{10pt}

% ------------------------------- Abstracts/Short Papers -------------------------------
\global\cvbibfirsttrue
\renewcommand{\cvbiblabel}{Abstracts/Short Papers}
\renewcommand{\cvbibtag}{}
\printbibliography[keyword=abstracts, env=cvpub, heading=none]

\vspace{10pt}

% ------------------------------- Talks -------------------------------
\global\cvbibfirsttrue
\renewcommand{\cvbiblabel}{Talks}
\renewcommand{\cvbibtag}{}
\printbibliography[keyword=talks, env=cvpub, heading=none]

\vspace{10pt}

% ------------------------------- Professional Outreach -------------------------------
\cvrow{Professional Outreach}{Reviewer, \textit{Medical Image Analysis}.}{2026}
\cvrow{}{Secretary, VALIANT Attempt, Vanderbilt University.}{2024 -- Present}
\cvrow{}{Reviewer, \textit{Scientific Reports}.}{2025}
\cvrow{}{Reviewer, \textit{Journal of Medical Imaging}.}{2024}

\vspace{10pt}

% ------------------------------- Honors and Awards -------------------------------
\cvrow{Honors and Awards}{Third Place, 3-minute Poster Competition, SPIE Medical Imaging}{February 2026}
\cvrow{}{Robert F. Wagner All-Conference Best Student Paper Finalist, SPIE Medical Imaging}{February 2025}
\cvrow{}{Second Place, 3-minute Poster Competition, SPIE Medical Imaging}{February 2025}
\cvrow{}{Graduate School Travel Grant, Vanderbilt University}{November 2024}
\cvrow{}{University Graduate Fellowship, Vanderbilt University}{July 2024}
\cvrow{}{Thomas R. Armstrong Award of Excellence for Outstanding Electrical Engineering Achievement, University of Dayton Department of Electrical and Computer Engineering}{May 2023}
\cvrow{}{Outstanding Collaborative Pianist, University of Dayton Department of Music}{May 2023}
\cvrow{}{Dean's List, University of Dayton}{August 2019 -- May 2023}
\cvrow{}{SPIE Travel Grant}{February 2023}
\cvrow{}{Best Poster in Directorate, ORNL Summer Intern Symposium, Oak Ridge National Laboratory}{August 2022}
\cvrow{}{Honors Thesis Fellowship, University of Dayton}{April 2022}
\cvrow{}{Finalist, International Science and Engineering Fair}{May 2018}

\vspace{10pt}

% ------------------------------- Skills -------------------------------
\cvpar{Skills}{Deep learning for medical imaging and data imputation using Python/PyTorch}

\vspace{4pt}
\cvpar{}{Adept at statistical analysis and interpretation}

\vspace{4pt}
\cvpar{}{Technical writing and typesetting using Microsoft Office and LaTeX}

\vspace{4pt}
\cvpar{}{Effective scientific communication and visual aid preparation}

\vspace{4pt}
\cvpar{}{Strong critical thinking and analytical skills}

\end{document}
